\documentclass[11pt]{article}
\usepackage[margin=1in]{geometry}
\usepackage{amsmath,amssymb,amsthm}
\usepackage{boldtensors}
\usepackage{hyperref}
\usepackage{booktabs}

\newcommand{\dd}{\mathrm{d}}

\title{Design: Variational Transfer for Impedance-Overlap Schwarz \\
on Non-Conforming Meshes}
\author{Alejandro Mota}

\begin{document}
\maketitle

\section{Problem statement}

The impedance-overlap transmission condition
\begin{equation}
  ~t_k - ~t_j
  + ~Z\,\bigl(\dot{~u}_k - \dot{~u}_j\bigr)
  + \alpha\,\bigl(~u_k - ~u_j\bigr) = ~0
  \qquad \text{on } \Gamma_k
  \label{eq:tc}
\end{equation}
is now exactly conservative on conforming meshes (0.08\% PU-energy error
over 1\,ms with the consistent partner traction, matching the
exact-DBC reference).  On the non-conforming cantilever pair
($h = 8.47$ vs.\ $6.35$\,mm, 4:3) the net error also looks small
($\pm 1$\%), but the interface-work instrumentation shows that this is
a cancellation of two error channels, each of which moves roughly
\emph{half the system energy per millisecond}:

\begin{center}
\small
\begin{tabular}{lccc}
\toprule
1\,ms, II, PU metric & conforming & non-conforming $h$ & non-conforming $h/2$ \\
\midrule
dashpot (spurious) work & $-42.6$\,J \,(2.8\%) & $-606.5$\,J \,(40\%) & $-598.5$\,J \,(40\%) \\
traction-channel work & --- & $+677$\,J \,(45\%) & $+768$\,J \,(51\%) \\
net PU-energy error & $1.5$\% $\to$ 0.08\% & $-0.75$\% (growth) & $+1.5$\% (loss) \\
mean velocity jump (RMS) & 8.9\,m/s & 40\,m/s (growing 15$\to$60) & 27\,m/s \\
\bottomrule
\end{tabular}
\end{center}

Three facts drive the design.  (i) The channels do \emph{not} shrink
under mesh refinement (dashpot work unchanged at $h/2$), so resolution
cannot fix the interface; the pointwise-interpolation transfer feeds
the condition an $O(1)$ error on mesh-scale content.  (ii) The net
error is the difference of two large channels and even flips sign
between $h$ and $h/2$ --- it is not robust, which is the production
fragility.  (iii) The velocity jumps grow through the run: each
interface pass leaves mesh-scale residue that reverberates, a bounded
version of the same feedback that makes DBC overlap coupling
exponentially unstable.

The theoretical requirements are Gander--Halpern--Nataf's three
conditions for discrete energy convergence of waveform-relaxation
coupling on nonmatching grids: (i) Robin (not Dirichlet) transmission,
(ii) consistent boundary-cell discretization of the transmission
operator, (iii) an $L^2$-contractive intergrid transfer.  The
impedance-overlap BC satisfies (i); its weak (surface-mass-weighted)
assembly satisfies (ii); pointwise interpolation violates (iii).
Replacing the transfer with the variational ($L^2$-projection, known as
mortar in the domain-decomposition literature) transfer onto the
receiving trace space satisfies all three simultaneously.  Note that
weak-$L^2$ enforcement of \emph{Dirichlet} overlap coupling did not
remove the DBC instability --- consistent with the theory, since DBC
violates (i) regardless of the transfer; the combination proposed here
is the first that meets all three conditions.

\section{Operator inventory (existing machinery)}

All operators needed for the first stage already exist:

\begin{center}
\small
\begin{tabular}{lll}
\toprule
operator & code & status \\
\midrule
$~W_k = \int_{\Gamma_k} ~N_k ~N_k^\top \dd\Gamma$ &
  \texttt{get\_square\_projection\_matrix} &
  already on the BC \\[0.3em]
$~L_{kj} = \int_{\Gamma_k} ~N_k(~x)\, ~N_j^{\mathrm{vol}}(~x)^\top \dd\Gamma$ &
  \texttt{get\_overlap\_rectangular\_projection\_matrix} &
  exists (weak-DBC path); \\
  & & needs signature relaxed \\[0.3em]
$~P_{kj} = ~W_k^{-1} ~L_{kj}$ &
  weak-DBC \texttt{dirichlet\_projector} &
  same construction \\[0.3em]
surface-to-surface $~R_{kj}$ &
  \texttt{get\_rectangular\_projection\_matrix} &
  exists (nonoverlap path); \\
  & (\texttt{project\_point\_to\_side\_set}) & needed in stage 2 \\[0.3em]
partner one-sided weak flux &
  \texttt{build\_consistent\_traction\_patch!}, &
  exists (conforming); \\
  & \texttt{consistent\_partner\_traction} & generalized in stage 2 \\
\bottomrule
\end{tabular}
\end{center}

$~L_{kj}$ is assembled by quadrature over $\Omega_k$'s boundary facets;
at each quadrature point the partner \emph{volume} element is located
(\texttt{find\_point\_in\_mesh}) and its shape functions evaluated ---
exactly the variational cross-mass for an overlap boundary, which is
interior to the partner and therefore has no partner surface mesh.

Two structural properties come for free.  First,
$~P_{kj} = ~W_k^{-1} ~L_{kj}$ is the $L^2(\Gamma_k)$-orthogonal
(Galerkin) projection onto $\Omega_k$'s trace space, hence
non-expansive in $L^2$: GHN condition (iii) holds by construction,
for any quadrature, because it is a projection.  Second, partition of
unity of the partner shape functions gives
$~L_{kj} \mathbf{1} = ~W_k \mathbf{1}$ for \emph{any} quadrature rule,
so constants (rigid translations) transfer exactly regardless of
quadrature accuracy.

\section{Design}

\subsection{Stage 1: variational projection of the transferred fields}

Replace pointwise interpolation of the three partner fields with the
variational projection, leaving everything downstream of the nodal values
unchanged:
\begin{equation}
  \hat{\dot{~u}}_j = ~P_{kj}\, \dot{~u}_j, \qquad
  \hat{~u}_j = ~P_{kj}\, ~u_j, \qquad
  \hat{~\sigma}_j = ~P_{kj}\, ~\sigma_j^{\mathrm{rec}}
  \quad \text{(six Voigt components),}
  \label{eq:stage1}
\end{equation}
then form the traction $\hat{~t}_j = \hat{~\sigma}_j\, ~n_k$ and the
tensor-impedance term nodally exactly as the current code does.  This
is a drop-in change: \texttt{apply\_bc\_detail}'s per-node
interpolation loop becomes three (or nine, per component) small dense
matrix--vector products with a projector precomputed at setup.  The
P/S-split tensor $~Z(~n)$ continues to act at the receiving nodes, so
no re-derivation of the self terms or the tangent is needed; the
variational crime of nodal application is unchanged from the current
scheme and identical on both branches.

Implementation touch points:
\begin{itemize}
\item \textbf{YAML}: new BC option \texttt{transfer:} with values
      \texttt{pointwise} (current behavior, initial default) and
      \texttt{variational}.
\item \textbf{Types}: field \texttt{variational\_projector::Matrix\{Float64\}}
      ($n_k \times N_j$, dense; boundary node count times partner node
      count, small for realistic interfaces) on
      \texttt{SolidMechanicsImpedanceOverlapSchwarzBoundaryCondition}.
\item \textbf{Setup} (\texttt{compute\_impedance\_overlap\_schwarz\_projectors!}):
      relax \texttt{get\_overlap\_rectangular\_projection\_matrix}'s
      signature to accept the impedance BC (it only touches fields both
      BC types share), build $~L_{kj}$, factorize the existing
      $~W_k$ (Cholesky) once, store $~P_{kj}$.  Rebuilt on BC swaps,
      like the consistent-traction patch.
\item \textbf{Evaluation}: factor the partner-field sampling shared by
      \texttt{apply\_bc\_detail} and
      \texttt{report\_impedance\_interface\_work} into one helper that
      dispatches on the transfer mode, so the instrumentation always
      measures what the BC actually applies.
\item \textbf{Precedence}: on node-aligned interfaces $~L_{kj}$ has a
      single unit entry per row and $~P_{kj}$ reduces to the identity;
      the consistent-traction patch (\texttt{partner traction: auto})
      continues to take priority for the traction, and the kinematic
      projection becomes a no-op.  The conforming limit is therefore
      automatic.
\end{itemize}

Quadrature note: the integrand of $~L_{kj}$ is only piecewise smooth
on $\Gamma_k$ (partner elements change within a receiving facet), so
the default facet rule carries a consistency error.  Because
contractivity and constant-transfer hold for any rule, this error
degrades accuracy, not stability.  The cheap knob is an elevated or
subdivided facet rule for the $~L_{kj}$ assembly only (setup cost);
a true common-refinement (2D polygon overlay on the flat interface)
is the exact alternative and is deferred until measurements justify
it.

What stage 1 answers: how much of the 600\,J churn is the
\emph{transfer} (interpolation aliasing of mesh-scale content) versus
the \emph{recovered-stress traction}.  The conforming ladder bounds
what recovery-quality traction can achieve ($\approx$1.5\%-class); if
variational transfer restores $O(h)$ scaling of the channels, refinement
plus stage 2 closes the rest.

\subsection{Stage 2: consistent traction across the interface}

The conforming analysis showed the traction channel is the dominant
residual once transfer is clean, and that nodal recovery cannot
represent the traction carried by mesh-scale content.  The consistent
(one-sided partial assembly) traction requires a partner facet surface
bounding a one-sided element patch; on a non-conforming interface
$\Gamma_k$ cuts through partner elements, so the patch does not exist
on $\Gamma_k$ itself.  Two generalizations, in order of increasing
cost:

\paragraph{2a: offset facet surface + surface variational.}  Take
$\tilde\Gamma_j$ = the partner facet surface immediately exterior to
$\Gamma_k$ (the boundary between the first exterior element layer and
the rest; the patch-construction logic already classifies these
elements).  The consistent flux is exact on $\tilde\Gamma_j$ by the
same partial assembly as the conforming case, tested with partner
surface functions.  Transfer the resulting weak traction to
$\Gamma_k$'s trace space with the surface-to-surface variational transfer
($\texttt{get\_rectangular\_projection\_matrix}$ generalized to accept
the internal surface in place of a side set, or an equivalent
assembled operator), i.e.\
$~W_k \hat{~t}_j = ~R_{k\tilde\jmath}\, ~W_{\tilde\jmath}^{-1}
\lambda_{\tilde\jmath}$.
The surface offset introduces an $O(h_j)$ modeling error
($\tilde\Gamma_j$ is within one partner element of $\Gamma_k$), but it
acts on the \emph{exact discrete flux} rather than on recovered
stress; whether that trade wins is an experiment, not a theorem ---
measure before committing further.

\paragraph{2b: cut-cell partial assembly (flat interfaces).}  For a
planar $\Gamma_k$ (the common production configuration), integrate the
partner's momentum residual over the plane-clipped parts of the cut
elements, testing with $\Omega_k$'s surface functions extended
constant along the interface normal.  This reproduces the
conforming construction exactly in the limit and removes the offset
error of 2a, at the cost of new machinery (hexahedron--half-space
clipping and quadrature on the clipped polyhedra).  Deferred unless 2a
measures short of the target.

\subsection{Validation and targets}

The metric is the instrumentation's channel decomposition, not the net
energy: success is driving the churn itself down, since the current
scheme already fakes a small net by cancellation.

\begin{enumerate}
\item \textbf{Conforming regression}: variational transfer on the
      node-aligned pair must reproduce the pointwise results exactly
      ($~P = ~I$); the consistent-traction 0.08\% result must be
      unaffected.  Extend test 103/104 asserts with the projector
      identity check.
\item \textbf{Non-conforming, stage 1}: dashpot work
      $\ll 600$\,J at $h$, and restored $O(h)$ decay at $h/2$;
      velocity-jump growth curve flattened.  Expected end state
      $\approx$ conforming recovered-stress class (few percent), now
      robust rather than cancellation-balanced.
\item \textbf{Non-conforming, stage 2}: traction channel reduced
      toward the consistent-traction class; net PU error
      $\lesssim 1$\% with both channels individually small; sign and
      magnitude stable under refinement.
\item \textbf{Stability margin}: rerun the impedance-scale sweep ---
      with contractive transfer the growth branch above scale 1 should
      soften, and the DBC-limit blow-up should set in later; this is
      the direct measure of production robustness.
\item \textbf{Cost}: setup $O(n_{\mathrm{facets}} \times n_{\mathrm{qp}})$
      point locations (same as the weak-DBC path); per-iteration cost
      \emph{decreases} (dense mat-vec replaces per-node interpolation
      loops; no recovery solve once stage 2 lands).
\end{enumerate}

\section{Stage 1 measurements (2026-07-03)}

Stage 1 was implemented (\texttt{transfer: variational}, commit-level
detail in the repository) and measured on the instrumented 1\,ms
non-conforming benchmark:

\begin{center}
\small
\begin{tabular}{lcccc}
\toprule
1\,ms, II, PU metric & pointwise $h$ & pointwise $h/2$ &
  variational $h$ & variational $h/2$ \\
\midrule
dashpot work [J] & $-606.5$ & $-598.5$ & $-726.1$ & $-483.9$ \\
traction-channel work [J] & $+677$ & $+768$ & $+880$ & $+700$ \\
net PU-energy error & $-0.75$\% & $+1.5$\% & $+6.2$\% (growth) & $+0.7$\% (growth) \\
mean velocity jump [m/s] & 40.0 & 27.0 & 37.8 & 28.0 \\
\bottomrule
\end{tabular}
\end{center}

Three conclusions.  First, the variational transfer restores an
$h$-dependence that pointwise interpolation lacks (dashpot churn
$726 \to 484$\,J under refinement, versus unchanged), but the decay is
slow and the channel magnitudes remain in the hundreds of joules ---
the error class is unchanged.  Second, at the coarse resolution the
net error \emph{worsens} (6.2\% growth versus 0.75\%): under pointwise
interpolation the kinematic-transfer error and the recovered-traction
error come from the same interpolant and partially cancel inside the
transmission condition; projecting the kinematics decorrelates them,
and the interface becomes net-injecting.  Third, the traction channel
is transfer-dependent and did not improve ($+880$\,J at $h$): the
recovered-stress traction remains the dominant polluter.

The design consequence: \textbf{stage 2 is not optional}.  The
kinematic variational transfer should land \emph{together} with a
variational (consistent) traction --- shipping stage 1 alone is
counterproductive at production resolutions.

\paragraph{Quadrature experiment.}  Elevating the facet quadrature of
$~L_{kj}$ (a $4\times4$ sub-cell rule via the new
\texttt{transfer quadrature subdivisions} option, sub-cells of
2.1\,mm against 6.35\,mm partner elements) separates the quadrature
consistency error from the traction-channel error at $h$: the net
error improves from $+6.2$\% to $+2.4$\% growth --- so the default
rule's consistency error was a material contributor to the stage-1 net
degradation and the subdivided rule should be used whenever the
facet-size ratio is coarse --- while the channel magnitudes grow
(dashpot $-726 \to -857$\,J, traction $+880 \to +941$\,J): the more
faithful projection captures more of the partner's mesh-scale content,
which the dashpot must then dissipate.  Accurate transfer quadrature
therefore does not change the conclusion: the recovered-stress
traction channel dominates, and stage 2 remains the required fix.

\section{Stage 2a measurements (2026-07-03)}

Stage 2a was implemented (offset-surface consistent traction with the
surface-to-surface transfer $~T = ~R \tilde{~W}^{-1}$) and verified
against a manufactured uniform uniaxial-stress state: the transferred
weak traction reproduces $~W (\sigma ~n)$ to $10^{-11}$ relative error
on the non-conforming pair, confirming the construction.  (A uniform
uniaxial-\emph{strain} state fails the same test by design: its
lateral confinement tractions are correctly picked up by the one-sided
assembly on the patch's lateral boundary strips.)  The instrumented
benchmark, however, is decisively negative, completing a fidelity
ladder at $h$:

\begin{center}
\small
\begin{tabular}{lcc}
\toprule
1\,ms, II, PU metric, $h$ & dashpot work & net PU error \\
\midrule
pointwise + recovered stress & $-607$\,J & $-0.75$\% \\
variational + recovered stress & $-726$\,J & $+6.2$\% (growth) \\
variational, subdivided + recovered stress & $-857$\,J & $+2.4$\% (growth) \\
variational, subdivided + consistent traction (2a) & $-1129$\,J & $+9.3$\% (growth) \\
\bottomrule
\end{tabular}
\end{center}

Each fidelity increase makes the interface churn \emph{worse}.  The
reading: on non-matching grids the round trip cannot be the identity;
whatever content cannot be represented across the interface must be
dissipated, or it aliases and reverberates.  Nodal stress recovery had
been acting as a beneficial low-pass filter; the exact discrete flux
transmits the partner's full mesh-scale content into a trace space
that cannot carry it, and the feedback loop (growing velocity jumps)
is pumped harder.  Fidelity without a mechanism that removes
non-transferable content is anti-stabilizing.

A second structural defect is now visible: the traction arrives
through $~T$ (built on the offset surface) while the velocity arrives
through $~P$ (built from volume interpolation) --- two different
filters.  The impedance condition has characteristic structure,
$t_k - Z \dot u_k = -(t_j + Z \dot u_j)$ up to signs: for outgoing
waves the combination $g_j = ~t_j + ~Z \dot{~u}_j + \alpha ~u_j$ is
smooth even when $~t_j$ and $\dot{~u}_j$ separately carry mesh-scale
standing content, because that content cancels \emph{within} $g_j$.
Transferring the two pieces through different operators destroys the
cancellation at exactly the mesh scale where it matters.

\paragraph{Proposed stage 2c: single-operator characteristic
transfer.}  Assemble the partner's full Robin datum on the offset
surface --- $\tilde{~g} = -\tilde{\lambda} + \tilde{~W}\,
(~Z\,\dot{~u}_j + \alpha\,~u_j)\big|_{\tilde\Gamma}$, using the
partner's own nodal trace values on $\tilde\Gamma$ (no interpolation
at all) --- and transfer it with the single operator $~T$.  The
mesh-scale cancellation then happens on the partner side, before
transfer.  Implementation cost is small: keep $\tilde{~W}$ and the
$\tilde\Gamma$ node list in the patch, and bypass $~P$ in the
right-hand side.  This was implemented and measured; see the
following section --- the hypothesis did not survive.

\section{Stage 2c measurements and revised conclusions (2026-07-03)}

Stage 2c was implemented (single-operator characteristic transfer of
the complete Robin datum, assembled on $\tilde\Gamma$ from the
partner's own nodal trace values) and benchmarked identically.  The
result is statistically indistinguishable from stage 2a: dashpot work
$-1084$ vs $-1129$\,J, net $+9.8$ vs $+9.3$\% growth, the same
velocity-jump growth.  The characteristic-cancellation hypothesis is
refuted: the operator pairing between traction and velocity is not the
driver.

The complete five-variant ladder at $h$ now reads: pointwise +
recovered $-607$\,J / $-0.75$\%; variational + recovered $-726$\,J /
$+6.2$\%; + subdivided quadrature $-857$\,J / $+2.4$\%; + consistent
traction (2a) $-1129$\,J / $+9.3$\%; characteristic (2c) $-1084$\,J /
$+9.8$\%.  Every fidelity increase pumps the churn.  The conclusion is
structural, not implementational (both stages pass the uniform-stress
patch test to $10^{-11}$): \textbf{on non-matching grids, content that
cannot be represented across the interface must be dissipated, and
smoothing the transferred data is beneficial --- the recovered-stress
low-pass filter was doing load-bearing stabilization work.}  The
parameter-free pointwise + recovered-stress scheme (with the partner
traction term and consistent recovery) is empirically the best of the
five variants at production resolution.

\paragraph{Where this leaves the design.}
(i) For production non-conforming runs, keep \texttt{transfer:
pointwise} with recovered stress as the default; the variational
machinery remains valuable for its conforming-limit exactness, the
quadrature knob, and as diagnostic infrastructure.
(ii) The principled next idea is a \emph{content-aware absorbing
term}: dissipate exactly the non-transferable component by penalizing
$(~I - \Pi)\,\dot{~u}_k$ with the characteristic impedance, where
$\Pi$ is the round-trip transfer operator --- absorb what cannot
cross, transmit what can.  This turns the accidental stabilization of
the recovery filter into a designed mechanism with a conforming limit
of zero.
(iii) A cheap discriminating diagnostic first: rerun the benchmark
with a dissipative integrator (Newmark $\gamma > 1/2$ or HHT-$\alpha$)
to confirm the churn is carried by the high-frequency band that
subdomain-level damping removes.

\section{Content-aware absorption and the closing diagnostic
(2026-07-03)}

The content-aware absorbing term was implemented (BC option
\texttt{content aware absorption}): an additional dissipative self
term $~W ~Z (~I - \Pi)\,\dot{~u}_k$ with its Newmark tangent, where
$\Pi$ is the $~W$-orthogonal projection onto the span of the
variational projector's columns --- the \emph{transferable} space at
the boundary.  The filter has exactly the designed structure: it
annihilates constants ($~P\mathbf{1} = \mathbf{1}$ puts rigid motions
in the span), vanishes identically on node-aligned interfaces and on
the coarse side of the study pair (everything a coarse trace holds is
representable by the fine partner), is nonzero only on the fine side
($\lVert ~F \rVert = 0.44$), and is provably dissipative
($v^\top ~W ~F\, v \ge 0$ per scalar channel; asserted in test 105).

\textbf{The measured effect is null}: the absorber drains $4.6$\,J
over 1\,ms while the channels and the net remain identical to the
pointwise baseline.  Because the operator is verified, the null is
informative: \emph{the reverberating content lives inside the
transferable space}.  Both grids can represent the mesh-scale content
that churns; they disagree about it because their discrete dispersion
relations differ.  No transfer-side construction can remove a
dispersion disagreement, by construction --- which retroactively
explains the entire fidelity ladder.

The closing diagnostic confirms the mechanism: rerunning the pointwise
baseline with a dissipative integrator (Newmark $\gamma = 0.6$,
$\beta = 0.3025$) collapses the churn from $-607$ to $-231$\,J and,
decisively, the velocity jumps \emph{decay} ($11.8 \to 5.6$\,m/s
between run halves) instead of growing ($33 \to 47$): the
reverberation loop is dead.  The churn is carried by the
high-frequency band that subdomain-level damping removes.  ($\gamma =
0.6$ is a blunt instrument --- it also removed 87\% of the physical
energy on this broadband benchmark; production should use mild
HHT-$\alpha$-class damping concentrated near the grid Nyquist.)

\paragraph{Final recommendations.}
(i) \textbf{Conforming interfaces}: consistent traction ---
0.08\%, exact-DBC class.
(ii) \textbf{Non-conforming production runs}: pointwise transfer +
recovered-stress traction with the physical impedance (parameter-free),
paired with mild high-frequency integrator damping; quantify the
damping trade-off with a dissipation sweep.
(iii) The variational machinery (projector, subdivided quadrature,
offset-surface flux, characteristic transfer, content filter) is
retained: it provides the conforming-limit exactness, the quadrature
control, the instrumentation-grade diagnostics that produced these
conclusions, and patch-test-verified building blocks for future
transmission-condition work.
(iv) The OSM-level lesson: the non-conforming overlap interface error
is dominated by the discrete dispersion disagreement of high-frequency
content between the subdomain discretizations --- a property of the
grids, not of the transmission condition.  The impedance dashpot is
the correct interface-level treatment and is already in place; the
remaining energy budget belongs to the time integrator.

\section{Summary}

Stage 1 is a low-risk, high-information change that reuses the square
($~W_k$) and rectangular ($~L_{kj}$) projection machinery verbatim:
variationally project $\dot{~u}_j$, $~u_j$, and $~\sigma_j^{\mathrm{rec}}$
onto the receiving trace space and keep the rest of the impedance
pipeline unchanged.  The conforming limit and the consistent-traction
path are preserved automatically.  The stage-1 measurements show that
contractive kinematics alone do not reduce the interface churn ---
they restore $h$-decay but break the error cancellation that kept the
pointwise scheme's net small --- so the variational transfer and the
stage-2 variational traction must ship together, with the channel
decomposition (not the net energy) as the acceptance metric.

\end{document}
